\chapter{State $s$ and Action $a$}
\label{ch:states_actions}

\begin{studentq}
``Action `a' is essentially state `s' because a is the causal factor of s. Why then differentiate these two? A new move a means a new position s. Isn't it?''
\end{studentq}

You are right that in chess they are linked deterministically:
\begin{equation}
  s' = \text{Transition}(s, a) \quad \text{(play move } a \text{ in position } s, \text{ arrive at } s'\text{)}
\end{equation}

So why keep two symbols? Because they live in different places in the tree, and they hold different numbers.

\begin{center}
% fig_c_circles_arrows.tex -- State s (Circles) vs Action a (Arrows)
% Annotation boxes sit at explicit staggered coordinates rather than on the
% edges: the four edges converge at the root, so pos=... placement collided and
% truncated the labels. Outer pair on the upper row, inner pair on the lower.
\begin{tikzpicture}[scale=0.92,
    poscircle/.style={circle, draw=NavyBlue, fill=PaleGrey, line width=1.1pt,
                      minimum size=23mm, align=center, font=\scriptsize},
    childcircle/.style={circle, draw=WorkGrey!70, fill=white, line width=0.8pt,
                        minimum size=17mm, align=center, font=\tiny},
    movearrow/.style={-{Stealth[length=2.8mm]}, draw=BuildBlue, line width=1.3pt},
    annotbox/.style={draw=DataTeal, fill=DataTeal!8, rounded corners=2pt,
                     font=\tiny, inner sep=2pt, align=center},
    dimbox/.style={draw=WorkGrey!60, fill=WorkGrey!5, rounded corners=2pt,
                   font=\tiny, inner sep=2pt, align=center}
  ]

  % Root circle (State s)
  \node[poscircle] (root) at (0, 3.4) {%
    \textbf{State $s$ (Root)}\\[0.3ex]
    White to move\\
    \color{NavyBlue}\textbf{$V(s) = +0.97602$}%
  };

  % Annotation for Root
  \node[draw=NavyBlue, fill=NavyBlue!10, rounded corners=2pt, font=\scriptsize,
        align=left, anchor=west] at (2.6, 3.4) {%
    \textbf{CIRCLES = POSITIONS (States $s$)}\\
    The board setup: 64 squares, side to move.\\
    \textbf{Single-pass evaluation $V(s)$ lives here.}%
  };

  % Child circles (States s')
  \node[childcircle] (c1) at (-5.1, -1.9) {\textbf{State $s_1'$}\\(after Kd6)\\$V = +0.96766$};
  \node[childcircle] (c2) at (-1.7, -1.9) {\textbf{State $s_2'$}\\(after Kf6)\\$V = +0.98598$};
  \node[childcircle] (c3) at (1.7, -1.9) {\textbf{State $s_3'$}\\(after Kf5)\\$V = \text{unvisited}$};
  \node[childcircle] (c4) at (5.1, -1.9) {\textbf{State $s_4'$}\\(after Kd5)\\$V = \text{unvisited}$};

  % Arrows (Actions a) -- drawn first so the labels sit on top of them
  \draw[movearrow] (root) -- (c1);
  \draw[movearrow] (root) -- (c2);
  \draw[movearrow, draw=WorkGrey!50] (root) -- (c3);
  \draw[movearrow, draw=WorkGrey!50] (root) -- (c4);

  % Action labels: outer pair high, inner pair low, so nothing collides
  \node[annotbox] at (-4.05, 1.35)
    {\textbf{Action $a_1$: Kd6}\\$P=45.13\%$\\$n=1,\ Q=+0.96766$};
  \node[annotbox] at (-1.55, -0.15)
    {\textbf{Action $a_2$: Kf6}\\$P=44.23\%$\\$n=1,\ Q=+0.98598$};
  \node[dimbox] at (1.55, -0.15)
    {\textbf{Action $a_3$: Kf5}\\$P=5.38\%$\\$n=0,\ Q=Q_{\mathrm{FPU}}$};
  \node[dimbox] at (4.05, 1.35)
    {\textbf{Action $a_4$: Kd5}\\$P=5.26\%$\\$n=0,\ Q=Q_{\mathrm{FPU}}$};

  % Annotation for Arrows
  \node[draw=BuildBlue, fill=BuildBlue!10, rounded corners=2pt, font=\scriptsize,
        align=center] at (0, -3.7) {%
    \textbf{ARROWS = MOVES (Actions $a$)}\\
    A move is a transition, not a position.\\
    \textbf{All search statistics ($P, n, W, Q$) live on the arrows!}%
  };
\end{tikzpicture}

\captionof{figure}{\textbf{The core ontological distinction.} Circles are board positions ($s$) where $V(s)$ lives; arrows are moves ($a$) where search statistics ($P, n, W, Q$) live.}
\label{fig:circles_arrows}
\end{center}

Read it out loud once:
\begin{itemize}[leftmargin=2em]
  \item \textbf{$s$ is a circle.} A board. Something you could set up with real pieces.
  \item \textbf{$a$ is an arrow.} A choice. Not a board --- you cannot set up an arrow.
\end{itemize}

\begin{established}[The Statistics Live on the Arrows, Not on the Circles]
The visit count $n$, the running total $W$, and the average quality $Q$ are all properties of an \textbf{edge}. When the guide writes $Q(s, a)$, it means: ``the quality of the arrow $a$ leaving circle $s$.'' It needs both symbols because the same move name (\texttt{Kd6}) means something different from a different position.
\end{established}

\begin{keyidea}[title={Your Own Summary, Which Was Correct}]
\textit{``$a$ is never a position. $s$ is the position prior to the moves $a$.''} Exactly. And the child positions $s'$ are circles again, each with their own four arrows leaving them. That is how the tree grows.
\end{keyidea}
